\section{Conclusion}

We introduced \phenaki, a model which is capable of generating variable length videos conditioned on a sequence of open domain text prompts. \phenaki uses \cvivit as video encoder. \cvivit is a new model which provides temporal-spatial compression while being auto-regressive in time.
The \cvivit model is a crucial part of \phenaki that allows it to generate variable length videos. We demonstrate how joint training on images and videos can improve the generation quality, and diversity, given the existence of much larger image-text dataset with order of magnitude more samples. The \phenaki model achieves good performance on video prediction, it can be used as to generate long videos conditioned on a text prompt. Additionally it is able to condition on both text and a starting frame. Finally, \phenaki is not limited to generating a video depicting a single concept or caption. It is actually able to generate longer coherent video stories based on a sequence of text prompts. The more complex narratives it can visualize demonstrate how this can become a great creative tool for story telling.

\section*{Ethics Statement}
While we have not explored potential downstream applications of the generative models described in this work, we believe Phenaki can have a positive impact in a variety of creative settings. In general, many of the samples from the model will not perfectly correspond to the input caption or the user’s intent; however, the end-user is likely to gain considerable time savings even if only one of the generated samples aligns with their intent. We thus foresee Phenaki being useful in eventually empowering users to accelerate their creativity, especially since the model can so quickly generate videos. Phenaki and similar models will be part of an ever-broad toolset for artists and non-artists alike, providing new and exciting ways to express creativity.

The flip-side of this acceleration and ease-of-use is the potential for harmful impact, as with many of the prior or concurrent work in generative modeling. An easy-to-use system like Phenaki can be repurposed for generating maliciously fake content and enable spreading of such content much easier. While the quality of the videos generated by Phenaki is not yet indistinguishable from real videos, getting to that bar for a specific set of samples is within the realm of possibility, even today. This can be particularly harmful if Phenaki is to be used to generate videos of someone without their consent and knowledge.

Like DALLE-2~\cite{ramesh2022hierarchical}, Imagen~\cite{saharia2022photorealistic}, Parti~\cite{PARTI} and others, Phenaki is trained on a collection of datasets that is known to encode a number of undesirable biases. LAION-400M~\cite{schuhmann2021laion} specifically has a variety of issues regarding violence, pornography, gore. While our primary image and video datasets have minimal traits like this, we did incorporate LAION-400M into our training and observed better results. In a currently training version of Phenaki, we use a set of datasets that minimizes such problems.

Taken together, these issues contribute to our decision not to release the underlying models, code, data or interactive demo at this time. Before we can do that, we want to focus our efforts on better understanding of data, prompt and output filtering. We would also like to more explicitly measure the biases encoded in the outputs of Phenaki, so that we can further mitigate them actively, either in the data, models or pre/post-processing steps.

